
We investigated entropy-based selective prediction on TriviaQA and discovered that model capacity acts as a binary gate for the validity of correlation-based statistical tests in uncertainty quantification research. Our empirical findings establish three validated contributions.

First, we confirm that token probability distributions from frozen LLM forward passes are fully extractable on factual QA tasks, with 100\% success rate across 500 TriviaQA examples. This validates extraction feasibility as technically unblocked. Second, we demonstrate that high max-probability, high-entropy disagreement cases exist in non-trivial proportions (8.20\% of predictions), confirming that entropy and max-probability capture different uncertainty dimensions and that the quadrant analysis framework is statistically sound. Third, we expose a critical experimental design requirement: models must provide non-zero variance in correctness for correlation-based validation; GPT-2 (117M parameters) produced 0\% accuracy, rendering Spearman correlation mathematically undefined (NaN). Based on scaling laws literature \citep{Roberts2020}, we hypothesize models below approximately 7B parameters invalidate such experiments on knowledge-intensive tasks, though we tested only GPT-2 and empirical confirmation across scales remains future work.

The third contribution is methodological rather than algorithmic. We do not claim entropy outperforms max-probability---our hypothesis remains untested due to invalid experimental conditions. Instead, we identify a failure mode that existing experimental practices overlook: zero-variance outcomes that make correlation tests undefined (NaN), not negative. This distinction matters because undefined tests indicate invalid experiments, while negative tests indicate unsupported hypotheses. Conflating the two leads to misinterpretation of methodological failures as substantive findings.

Our work establishes a practical guideline for future selective prediction research: verify that evaluation metrics exhibit non-zero variance before conducting correlation-based analyses. For factual QA specifically, this translates to a minimum model capacity threshold of approximately 7 billion parameters (extrapolated from Roberts et al., 2020, not empirically confirmed by our single-model study), or empirical confirmation that baseline accuracy exceeds 5\%. This simple pre-flight check---computing variance(correctness) and confirming it is non-zero---takes seconds and prevents wasted compute on invalid experimental runs.

The broader lesson extends beyond selective prediction. Correlation-based validation is ubiquitous in machine learning uncertainty quantification---calibration studies, ensemble comparison, conformal prediction evaluation. All share the same dependency: both variables must vary for correlation to be defined. Our framework for distinguishing infrastructure validation from hypothesis validation applies to any method where technical feasibility (can we extract signals?) is separate from performance claims (do signals predict quality?). Validating infrastructure first, before testing comparative performance, allows incremental progress even when hypotheses remain untested.

Looking forward, immediate next steps include re-running our experiment with Llama-2-7B as originally specified to determine whether entropy-based rejection outperforms max-probability under valid test conditions. If the hypothesis is supported, mechanism-testing sub-hypotheses can proceed to explain why entropy works (multi-modal uncertainty capture, top-5 entropy patterns, threshold-based rejection curves). If the hypothesis is refuted, alternative uncertainty measures (semantic entropy, ensemble disagreement) can be explored using the validated infrastructure we established.

Longer-term research directions include establishing model capacity thresholds across task types (reasoning, generation, classification) and developing automated experiment validity checkers that flag zero-variance outcomes before expensive compute is committed. Such tools would encode the methodological insight we discovered---that statistical tests have prerequisites, and violating them invalidates experiments regardless of implementation correctness.

We learned that infrastructure validation stands independent of hypothesis testing, that model capacity is not merely a performance variable but an experimental prerequisite, and that distinguishing between ''the test failed'' and ''the test could not run'' is essential for interpreting experimental outcomes correctly. Our infrastructure validation remains a contribution; our hypothesis awaits proper test conditions.
